\documentclass[10pt,a4paper]{amsart}
\usepackage[margin=19mm]{geometry}
\usepackage{amsmath,amssymb,mathtools}
\usepackage[scr=rsfs,cal=euler]{mathalfa}
\usepackage{xfrac}
\usepackage[unicode,hidelinks,bookmarks=false]{hyperref}
\newcommand{\ie}{i.e.\ }
\newcommand{\cf}{cf.\ }
\newcommand{\resp}{respectively\ }
\newcommand{\Subsection}{\subsection}
\providecommand{\nobreakdash}{\nobreak\hbox{-}}
\setlength{\emergencystretch}{2em}
\multlinegap0pt
\usepackage{bm}
\usepackage{bbm}
\usepackage{dsfont}
\usepackage{xspace}
\usepackage{cancel}
\usepackage{mathtools}
\usepackage{amsthm}
\makeatletter
\@ifundefined{theorem}{\newtheorem{theorem}{Theorem}}{}
\@ifundefined{definition}{\newtheorem{definition}[theorem]{Definition}}{}
\makeatother
\newcommand{\SN}{\mathfrak{S}}
\title[Reselection in the game of best choice]{Reselection in the game of best choice}
\author{Dillon Hanson, Brant Jones, Hyejin Kim}
\date{Source: arXiv:2607.27033v1 (CC BY 4.0)}
\begin{document}
\maketitle

\noindent\emph{Source and licence.} Reselection in the game of best choice. Version arXiv:2607.27033v1 (CC BY 4.0), distributed under the Creative Commons Attribution 4.0 International licence, as recorded in the arXiv licence metadata for this exact version and shown on the abstract page https://arxiv.org/abs/2607.27033v1. Full rights notice: licenses/arxiv-2607.27033-CC-BY-4.0.txt.\par\smallskip

\noindent\textit{Editorial scope.} Counted unit: the Theorem whose source label is \texttt{t:rnap} in the article (source0.tex lines 1281--1284, source label \texttt{t:rnap}), delivered together with the article's own complete original proof of it (source0.tex lines 1285--1338, one \texttt{proof} environment of 54 lines). No mathematics is written, repaired or re-derived: statement and proof are the article's own text line for line. Required context retained uncounted: d:sint (lines 1258--1273). Every definition, display and label of those lines is retained unchanged. Omitted from this package and not redistributed: 1--1257, 1274--1280, 1339--2494. Nothing inside a retained excerpt is omitted or altered: every retained body is delivered exactly as the article's lines are written, so no omission receipt of any kind accompanies this package. Citations: the retained excerpts carry no citation command at all, so no bibliography entry of the article is transcribed and no reference is redistributed. Reference adaptation: none was needed and none was made; every reference inside the delivered unit resolves inside it, so no unresolved reference or citation is produced. Printed slips kept literal: no printed word, symbol, hypothesis or number of the counted statement or of its proof was altered. Layout and class: the article's own class is \texttt{amsart} with packages including verbatim, amssymb, amsthm, amsmath, color, ifpdf, epsfig, lscape, hyperref, graphicx, tikz, ifthen; the wrapper is the lane's pinned \texttt{amsart} editorial wrapper at 10pt on A4, which restates the theorem-like environments the retained text needs and adds the article's own macro definitions that the retained text needs (\texttt{\textbackslash SN}), with extra wrapper packages (bm, bbm, dsfont, xspace, cancel, mathtools). The delivered numbering is the unit's own. Assets: no retained span contains a figure environment, a graphics-inclusion command, a PDF-page inclusion, a table or a TikZ picture; no image file is copied into this package, the package declares no asset dependency and no image file is redistributed with it. Licence: version arXiv:2607.27033v1 of this article is distributed under the Creative Commons Attribution 4.0 International licence, as recorded in the arXiv licence metadata for this exact version and shown on the abstract page https://arxiv.org/abs/2607.27033v1. A full rights notice is delivered at licenses/arxiv-2607.27033-CC-BY-4.0.txt. Characters: every retained excerpt is pure ASCII, so the delivered engine, XeLaTeX with its default Latin Modern text font, reports no missing character.
\begin{definition}\label{d:sint}
Suppose that $p$ is a prefix of size $k$.  Define an {\bf immediate child} of
$p$ to be a prefix $q$ of size $k+1$ that has its first $k$ entries in the same
relative order as $p$.  Each prefix $p$ has $k+1$ immediate children
distinguished by the value in their last position (as the complementary values
must all appear in the relative order specified by $p$).  Only one of the
immediate children ends in a left-to-right maximum (namely, the value $k+1$
itself); denote this child by $\bar p$.  

Then we define the {\bf interior payoff} function
\[ S^\circ(p) = \max\left( S(\bar p), S^\circ(\bar p) \right) + \sum\limits_{\text{other immediate children $q$ of $p$}} S^\circ(q) \]
recursively in terms of $S$ and $S^\circ$ values of prefixes having strictly greater size.
For initial conditions, we use that $S^\circ(p) = 0$ whenever $p$ has size $N$.  (We use the ``topological interior'' notation here to emphasize that the
interior payoff is a sum of certain strike payoffs from the interior of the
``prefix tree'' rooted at $p$, but never includes $S(p)$ itself.)
\end{definition}

\begin{theorem}\label{t:rnap}
The strategy that rejects each negative prefix and accepts the first positive
prefix encountered in the game of best choice is optimal.
\end{theorem}

\begin{proof}
Consider the tree of possible moves for a game of best choice.  Each node
consists of a prefix that we may as well assume ends in a left-to-right maximum
(for otherwise, the strike payoff is zero so the prefix is negative) or else
consists of a complete permutation of size $N$.  These nodes represent the
possible prefix flattenings that are presented to the player throughout the
game.

The nodes are partially ordered by containment of prefixes.  For example, the
root node of the game is the prefix $1$.  This node is directly connected to
prefixes $12$, $213$, $3124$, and others; the prefix $213$ is directly
connected to $2134$ as well as $31425$ and others.  At each node, the player
may accept the current candidate or reject the current candidate.  If no
prefixes are accepted, the player eventually encounters the permutations of
size $N$ and must accept the last candidate (whether they are best or not).

To determine the optimal strategy, begin at the permutations of size $N$.  We
view these as prefixes $p$ that are encountered by a player who has not yet
accepted any candidate.  Then, $S(p)$ is clearly $0$ or $1$ (depending on
whether the current candidate is $N$) while $S^\circ(p)$ is $0$ because there
is no way to reject a prefix that is already size $N$.  Thus, these prefixes
are all positive and accepting any of them is optimal (in fact, it is the
only allowable move left in the game).

Next, we argue by induction.  Suppose that our ``reject negativity, accept
positivity'' strategy is optimal when applied to any prefixes of size $k+1,
k+2, \ldots, N$.  Consider a prefix of size $k < N$.  Such a prefix necessarily
ends in a left-to-right maximum.

Let $\pi \in \SN_N$ represent the true sequence of candidate ranks in the
particular game being presented to the player and consider the moves available
in response to the prefix flattening $p$.  Our probabilities are conditioned on
$\pi \in \SN_N$ having $p$ as a prefix flattening.  By definition, accepting
the current candidate has a probability of winning given by $S(p)$ (divided by
a sum of $R_{J(\pi)}(x)$ terms over all $\pi$ that have $p$ as a prefix
flattening).  On the other hand, by Definition~\ref{d:sint}, rejecting the
current candidate has a probability of winning given by $S^\circ(p)$ (divided
by the same sum of $R_{J(\pi)}(x)$ terms over all $\pi$ that have $p$ as a
prefix flattening).  That is, the formula in Definition~\ref{d:sint} describes
the optimal win probabilities on each subtree from the immediate children; as
these subtrees are disjoint events whose union is the event we are calculating,
we simply sum the results.  We note that determining $S^\circ(p)$ explicitly
does require knowing the optimal strategy from this point in the game onward,
but we may assume that we have already computed $S^\circ(q)$ values and
assigned ``positive'' and ``negative'' status for each prefix $q$ with length
greater than $k$, by the induction hypothesis.  Thus, $S^\circ(p)$ is
well-defined.

Hence, the same ``reject negativity, accept positivity'' strategy is seen to be
optimal for the prefix $p$.  By induction, this strategy is optimal for all the
nodes in the tree.  In particular, $S^\circ(\emptyset) := \max(S(1),
S^\circ(1))$ is the probability of winning the game overall, playing from the
initial prefix flattening $1$ under the optimal strategy.
\end{proof}

\end{document}
